% GENERATED FILE. Edit canonical lessons, then run npm run generate:books.
% canonical-source-hash: fnv1a64:c287a247c22b9b2b
% canonical-lessons: BN-C454-songkoch, BN-C454-dvidha, BN-C454-aroshto, BN-C454-jorosoro, BN-C454-itostoto-kora

\chapter{Shyness and Hesitation}
\label{ch:bn-shyness-and-hesitation}
\hlchaptermodality{454}

\begin{chapteropening}
\textbf{By the end of this chapter:} \emph{I can say that I felt shy, awkward or unsure.}

Complete the last lesson of chapter 454: I can say that I felt shy, awkward or unsure.
\end{chapteropening}

\section[\texorpdfstring{\bn{সংকোচ} (sôngkoch)}{সংকোচ (sôngkoch)}]{\bn{সংকোচ} (sôngkoch) — reluctance; bashful hesitation}

\label{lesson:BN-C454-songkoch}



\begin{quote}
Before the new one: say the Bengali for to stand up to, then the Bengali for to back down.
\end{quote}

\subsection*{You'll want to know: \bn{সংকোচ}}
\textbf{\bn{সংকোচ}} — \emph{sôngkoch} — \textquotedblleft{}reluctance; bashful hesitation\textquotedblright{}.

\begin{quote}
\bn{সংকোচ} \bn{না} \bn{করে} \bn{তোমার} \bn{প্রশ্নটা} \bn{জিজ্ঞাসা} \bn{করো।}
\end{quote}

\begin{quote}
\emph{sôngkoch nā kôre tomār prôshnôṭā jigyāsā kôro.}
\end{quote}

\begin{quote}
Don't hold back: ask your question.
\end{quote}

\bn{সংকোচ} is the shrinking-back of a guest or a newcomer who does not want to impose: \bn{সংকোচ} \bn{করবেন} \bn{না}, please don't hold back.

\begin{cousinweb}
\textbf{\bn{সংকোচ}} (\emph{sôngkoch}) is a tatsama, Sanskrit \emph{saṃkoca}, \textquotedblleft{}contraction, shrinking together\textquotedblright{} (Monier-Williams): shrinking back.
\end{cousinweb}

\begin{grammarlens}[title={What you've built}]
One more word for describing what you lived and felt.
\end{grammarlens}

\subsection*{Your turn}
\begin{itemize}
\raggedright
  \item \emph{Say it:} \emph{sôngkoch}
  \item \emph{Say it:} \emph{sôngkoch}, once more
  \item \emph{Say it:} \emph{sôngkoch nā kôre tomār prôshnôṭā jigyāsā kôro.} — Don't hold back: ask your question.
  \item \emph{Recall:} say the Bengali for to stand up to, then the Bengali for to back down, then say \emph{sôngkoch} again
\end{itemize}

\begin{tcolorbox}[breakable,colback=teal!4,colframe=teal!35!black,title={Before you move on}]
What does \bn{সংকোচ} mean? (\textquotedblleft{}Reluctance; bashful hesitation\textquotedblright{}.) Say it once more.
\end{tcolorbox}
\section[\texorpdfstring{\bn{দ্বিধা} (dvidhā)}{দ্বিধা (dvidhā)}]{\bn{দ্বিধা} (dvidhā) — doubt; being in two minds}

\label{lesson:BN-C454-dvidha}



\begin{quote}
Before the new one: say the Bengali for to back down, then the Bengali for reluctance.
\end{quote}

\subsection*{You'll want to know: \bn{দ্বিধা}}
\textbf{\bn{দ্বিধা}} — \emph{dvidhā} — \textquotedblleft{}doubt; being in two minds\textquotedblright{}.

\begin{quote}
\bn{নতুন} \bn{চাকরি} \bn{নিয়ে} \bn{অরুণের} \bn{মনে} \bn{দ্বিধা} \bn{ছিল।}
\end{quote}

\begin{quote}
\emph{nôtun chākri niye ôruner mône dvidhā chhilô.}
\end{quote}

\begin{quote}
Arun was in two minds about the new job.
\end{quote}

\begin{cousinweb}
\textbf{\bn{দ্বিধা}} (\emph{dvidhā}) is a tatsama, Sanskrit \emph{dvidhā}, \textquotedblleft{}in two ways, divided\textquotedblright{} (Monier-Williams); Bengali turned the adverb into a noun, a mind split in two.
\end{cousinweb}

\begin{grammarlens}[title={What you've built}]
One more word for describing what you lived and felt.
\end{grammarlens}

\subsection*{Your turn}
\begin{itemize}
\raggedright
  \item \emph{Say it:} \emph{dvidhā}
  \item \emph{Say it:} \emph{dvidhā}, once more
  \item \emph{Say it:} \emph{nôtun chākri niye ôruner mône dvidhā chhilô.} — Arun was in two minds about the new job.
  \item \emph{Recall:} say the Bengali for to back down, then the Bengali for reluctance, then say \emph{dvidhā} again
\end{itemize}

\begin{tcolorbox}[breakable,colback=teal!4,colframe=teal!35!black,title={Before you move on}]
What does \bn{দ্বিধা} mean? (\textquotedblleft{}Doubt; being in two minds\textquotedblright{}.) Say it once more.
\end{tcolorbox}
\section[\texorpdfstring{\bn{আড়ষ্ট} (āṛôshṭô)}{আড়ষ্ট (āṛôshṭô)}]{\bn{আড়ষ্ট} (āṛôshṭô) — stiff; awkward and tongue-tied}

\label{lesson:BN-C454-aroshto}



\begin{quote}
Before the new one: say the Bengali for reluctance, then the Bengali for being in two minds.
\end{quote}

\subsection*{You'll want to know: \bn{আড়ষ্ট}}
\textbf{\bn{আড়ষ্ট}} — \emph{āṛôshṭô} — \textquotedblleft{}stiff; awkward and tongue-tied\textquotedblright{}.

\begin{quote}
\bn{নতুন} \bn{লোকের} \bn{সামনে} \bn{মীরা} \bn{আড়ষ্ট} \bn{হয়ে} \bn{যায়।}
\end{quote}

\begin{quote}
\emph{nôtun loker sāmne mīrā āṛôshṭô hôye jāy.}
\end{quote}

\begin{quote}
In front of new people Mira goes stiff and tongue-tied.
\end{quote}

\begin{cousinweb}
\textbf{\bn{আড়ষ্ট}} (\emph{āṛôshṭô}) is a Bengali word for stiffness, of the body or of the manner; its origin is not settled.
\end{cousinweb}

\begin{grammarlens}[title={What you've built}]
One more word for describing what you lived and felt.
\end{grammarlens}

\subsection*{Your turn}
\begin{itemize}
\raggedright
  \item \emph{Say it:} \emph{āṛôshṭô}
  \item \emph{Say it:} \emph{āṛôshṭô}, once more
  \item \emph{Say it:} \emph{nôtun loker sāmne mīrā āṛôshṭô hôye jāy.} — In front of new people Mira goes stiff and tongue-tied.
  \item \emph{Recall:} say the Bengali for reluctance, then the Bengali for being in two minds, then say \emph{āṛôshṭô} again
\end{itemize}

\begin{tcolorbox}[breakable,colback=teal!4,colframe=teal!35!black,title={Before you move on}]
What does \bn{আড়ষ্ট} mean? (\textquotedblleft{}Stiff; awkward and tongue-tied\textquotedblright{}.) Say it once more.
\end{tcolorbox}
\section[\texorpdfstring{\bn{জড়সড়} (jôṛôsôṛô)}{জড়সড় (jôṛôsôṛô)}]{\bn{জড়সড়} (jôṛôsôṛô) — huddled up; shrinking (with shyness or cold)}

\label{lesson:BN-C454-jorosoro}



\begin{quote}
Before the new one: say the Bengali for being in two minds, then the Bengali for stiff and awkward.
\end{quote}

\subsection*{You'll want to know: \bn{জড়সড়}}
\textbf{\bn{জড়সড়}} — \emph{jôṛôsôṛô} — \textquotedblleft{}huddled up; shrinking (with shyness or cold)\textquotedblright{}.

\begin{quote}
\bn{শীতে} \bn{কুকুরটা} \bn{জড়সড়} \bn{হয়ে} \bn{বসে} \bn{ছিল।}
\end{quote}

\begin{quote}
\emph{shīte kukurṭā jôṛôsôṛô hôye bôse chhilô.}
\end{quote}

\begin{quote}
The dog sat huddled up in the cold.
\end{quote}

\begin{cousinweb}
\textbf{\bn{জড়সড়}} (\emph{jôṛôsôṛô}) is a rhyming doubling whose first half probably belongs with Sanskrit \emph{jaḍa}, \textquotedblleft{}dull, numb with cold\textquotedblright{}, under which CDIAL lists Bengali forms (CDIAL 5090).
\end{cousinweb}

\begin{grammarlens}[title={What you've built}]
One more word for describing what you lived and felt.
\end{grammarlens}

\subsection*{Your turn}
\begin{itemize}
\raggedright
  \item \emph{Say it:} \emph{jôṛôsôṛô}
  \item \emph{Say it:} \emph{jôṛôsôṛô}, once more
  \item \emph{Say it:} \emph{shīte kukurṭā jôṛôsôṛô hôye bôse chhilô.} — The dog sat huddled up in the cold.
  \item \emph{Recall:} say the Bengali for being in two minds, then the Bengali for stiff and awkward, then say \emph{jôṛôsôṛô} again
\end{itemize}

\begin{tcolorbox}[breakable,colback=teal!4,colframe=teal!35!black,title={Before you move on}]
What does \bn{জড়সড়} mean? (\textquotedblleft{}Huddled up; shrinking (with shyness or cold)\textquotedblright{}.) Say it once more.
\end{tcolorbox}
\section[\texorpdfstring{\bn{ইতস্তত} \bn{করা} (itôstôtô kôrā)}{ইতস্তত করা (itôstôtô kôrā)}]{\bn{ইতস্তত} \bn{করা} (itôstôtô kôrā) — to hesitate}

\label{lesson:BN-C454-itostoto-kora}



\begin{quote}
Before the new one: say the Bengali for stiff and awkward, then the Bengali for huddled up.
\end{quote}

\subsection*{You'll want to know: \bn{ইতস্তত} \bn{করা}}
\textbf{\bn{ইতস্তত} \bn{করা}} — \emph{itôstôtô kôrā} — \textquotedblleft{}to hesitate\textquotedblright{}.

\begin{quote}
\bn{দরজার} \bn{সামনে} \bn{অরুণ} \bn{ইতস্তত} \bn{করল।}
\end{quote}

\begin{quote}
\emph{dôrjār sāmne ôrun itôstôtô kôrlô.}
\end{quote}

\begin{quote}
At the door Arun hesitated.
\end{quote}

Past (he/she …): \bn{তিনি} \bn{দরজায়} \bn{ইতস্তত} \bn{করলেন}, \emph{tini dôrjāy itôstôtô kôrlen}.

\begin{cousinweb}
\textbf{\bn{ইতস্তত} \bn{করা}} (\emph{itôstôtô kôrā}) is Sanskrit \emph{itastatas}, \textquotedblleft{}here and there, hither and thither\textquotedblright{} (Monier-Williams), with \emph{kôrā}, \textquotedblleft{}to do\textquotedblright{} (CDIAL 2814): going this way and that.
\end{cousinweb}

\begin{grammarlens}[title={What you've built}]
One more word for describing what you lived and felt.
\end{grammarlens}

\subsection*{Your turn}
\begin{itemize}
\raggedright
  \item \emph{Say it:} \emph{itôstôtô kôrā}
  \item \emph{Say it:} \emph{itôstôtô kôrā}, once more
  \item \emph{Say it:} \emph{dôrjār sāmne ôrun itôstôtô kôrlô.} — At the door Arun hesitated.
  \item \emph{Recall:} say the Bengali for stiff and awkward, then the Bengali for huddled up, then say \emph{itôstôtô kôrā} again
\end{itemize}

\begin{tcolorbox}[breakable,colback=teal!4,colframe=teal!35!black,title={Before you move on}]
What does \bn{ইতস্তত} \bn{করা} mean? (\textquotedblleft{}To hesitate\textquotedblright{}.) Say it once more.
\end{tcolorbox}
